 \documentclass{article}
\usepackage{tikz}
\usetikzlibrary{angles, quotes}
\usepackage{amssymb}
\usepackage{pgfplots}
\pgfplotsset{width=10cm,compat=1.9}
\usepgfplotslibrary{statistics}
\usepackage[left=2.54cm,right=2.54cm,top=2.54cm,bottom=2.54cm]{geometry}
\usepackage{tasks} %Paquete necesario para la producción de items horizontales para algunos ejercicios de tarea empleando el comando /begin{multicols}}{#}.../end{multicols} para ayudarnos a reducir las páginas
\usepackage{pdflscape} %comando que permite cambiar a orientación horizontal las hojas%
\usepackage{soul} %Paquete para permitir la compilación de acentos usando el comando {\hl{}}}\\
\sethlcolor{yellow(munsell)} %comando para definir el color para el subrayado usando el comando \hl
\usepackage[utf8]{inputenc}
\usepackage[spanish]{babel}
\usepackage{icomma}
\usepackage{ragged2e}
\usepackage{siunitx}
\usepackage{url}
\usepackage[colorlinks=true, urlcolor=blue,  linkcolor=black, citecolor=green]{hyperref}
\usepackage{pdfpages}
\usepackage{blindtext} %Paquete que genera texto ficticio (dummy text) con el comando: \blindtext
\usepackage{cancel} %in the preamble gives you four different modes of striking through: \cancel{text to cancel} draws a diagonal line (slash) through its argument, \bcancel{text to cancel} uses the negative slope (a backslash), \xcancel{text to cancel} draws an X (actually \cancel plus \bcancel), \cancelto{〈value〉}{〈expression〉} draws a diagonal arrow through the 〈expression〉pointing to the 〈value〉 (math-mode only)
\usepackage{csquotes}
\usepackage{afterpage}
\usepackage{parskip} 
\usepackage{float}
\usepackage{enumitem}
\usepackage{multicol}%Paquete que permite la creación aislada de columnas de texto. De esta forma se reduce la cantidad de páginas en nuestro documento
\usepackage{lipsum} %Paquete que genera texto de relleno al igual que \usepackage{blindtext}, se usa con el comando \lipsum[#-#] los #(hashtags) delimitan el número de páginas que deseamos ocupar del paquete lipsum para usarlos como relleno. 

% Definir el comando personalizado para la nota
\newcommand{\mynote}[1]{%
\begin{tikzpicture}[baseline=-0.75ex]
\node[draw, circle, fill=Apple Green!60, inner sep=2pt] (note) {\textbf{Nota:}};
\end{tikzpicture}%
\ \textit{#1}%
}

\newenvironment{Figura}
{\par\medskip\noindent\minipage{\linewidth}}
{\endminipage\par\medskip}
\usepackage{caption}
\usepackage[
backend=biber,
sorting=none,
url=true
]{biblatex} %bibliografía
\addbibresource{biblio.bib}
% Establecer el espacio entre las entradas de la bibliografía
\setlength{\bibitemsep}{1\baselineskip} % Puedes ajustar el valor según tus preferencias
\usepackage{amsmath,amsthm,amssymb,amsfonts}
\usepackage{pifont} %Permite la compilación del comando \xmark: es el símbolo de la tachita
\usepackage{mathtools} %permite la compilación de símbolos para matrices
\usepackage{empheq} %Paquete que se relaciona con \usepackage[most]{tcolorbox} para la creación de cuadros/cajas de colores para delimitar los resultados matemáticos o líneas de texto usando la línea de comando: \begin{empheq}[box={\mymath[colback=orange(webcolor)!70,drop lifted shadow]}]{equation*} \end{empheq}
\usepackage[most]{tcolorbox} %Paquete que se relaciona con \usepackage{empheq} para la creación de cuadros/cajas de colores para delimitar los resultados matemáticos o líneas de texto
\renewcommand{\qedsymbol}{$\blacksquare$}
\usetikzlibrary{positioning,decorations.pathreplacing} %Paquete que permite la compilación de llaves para cuadros sinópticos (brace diagram) 
\usepackage{schemata} %The schemata package is designed just for these brace diagrams
\usepackage{booktabs}

\newcommand\AB[2]{\schema{\schemabox{#1}}{\schemabox{#2}}} %comando o paquete necesario para crear cuadros sinópticos

\newtcbox{\mymath}[1][]{%
nobeforeafter, math upper, tcbox raise base,
enhanced, colframe=blue!30!black,
colback=blue!30, boxrule=1pt,
#1} %Comando importante para encasillar los resultados matemáticos en cajas de diferentes colores, se relaciona con los paquetes \usepackage{empheq} y \usepackage[most]{tcolorbox} 

\newenvironment{sysmatrix}[1]
{\left(\begin{array}{@{}#1@{}}}
{\end{array}\right)}
\newcommand{\ro}[1]{%
\xrightarrow{\mathmakebox[\rowidth]{#1}}%
}
\newlength{\rowidth}% row operation width
\AtBeginDocument{\setlength{\rowidth}{3em}}  %Comando importante para laproducción de lineas de operaciones entre matrices, método de Gauss o Gauss Jordan se relaciona con \newenvironment{sysmatrix}

%formato para cambiar el horario a español
\usepackage[spanish]{datetime2}
\DTMsetdatestyle{spanish}

\renewcommand{\today}{\DTMdisplaydate{\the\year}{\the\month}{\the\day }{-1}}
%%%%%%%%%%%%%%%%%%%%%%%%%%%%%%%%%%%%%%%%%

\usepackage{datetime}
\newcommand{\mycurrenttime}{\xxivtime}

%%%%%%%%%%%%%%%%%%%%%%%%%%%%%%%%%%%%%%%%%%%%%%%%%%%%%%%%%%%%%%%%%%%%%%%%%%%%%%%%%%%%%%%%%%%%%%%%%%%%%%%%%%%%%%%%%%%%%%%%%%%%%%%%%%%%%%%%%%%%%%%%%


%%%%%%%%%%%%%Caja de color para el título%%%%%%%%%%%%%%%%%%%%%%%%%%%%%%%%%%%%%%%%%%%%%%%%%%%%%%

\definecolor{myframecolor}{RGB}{1, 45, 75} % Prussian Blue
\definecolor{myboxcolor}{RGB}{0, 107, 92} % Apple Green

\newtcolorbox{mybox}{
enhanced,
colback=myboxcolor!7, % Color de fondo del cuadro
colframe=myframecolor, % Color del marco del cuadro
arc=0pt,
boxrule=1pt,
borderline west={2mm}{-10mm}{myframecolor}, % Borde en el lado izquierdo
sharp corners=southwest,
width=\linewidth,
before=\par\vspace{\bigskipamount}, % Espacio antes del cuadro
after=\par\vspace{\bigskipamount} % Espacio después del cuadro
}


\newcommand{\euler}{e} %Comando para producir letra e de euler
\newenvironment{remark}{\par\vfill\footnotesize % Vertical white space above the remark and smaller font size
\begin{list}{}{
	\leftmargin=80pt % Indentation on the left
	\rightmargin=60pt}\item\ignorespaces % Indentation on the right
\makebox[-2.5pt]{\begin{tikzpicture}[overlay]
		\node[draw=Horizon!60,line width=2.5pt,circle,fill=Horizon!25,font=\sffamily\bfseries,inner sep=4pt,outer sep=0pt] at (-19pt,5pt){\textcolor{Horizon}{Nota}};\end{tikzpicture}} % Blue Nota in a circle
\advance\baselineskip -1pt}{\end{list}\vskip5pt} % Tighter line spacing and white space after remark
\usepackage{graphicx}

\usepackage{titling}

\input{colores.tex}

%%Producir signos de cita textual``Asignatura''

\newcommand{\xmark}{\ding{55}} %%%comando que genera la tachita

\newenvironment{MyColorPar}[1]{%
\leavevmode\color{#1}\ignorespaces%
}{%
}%

%%%%%%%%%%%%%%%%%%% Título de la tarea, Nombre de alumno

\title{Clase \#32 - Electrodinámica}
\author{\emph{JA} $\boldsymbol{\mid}$}
\date{\today}

%%%%%%%%%%%%%%%%%%% Datos de la Materia

\usepackage{fancyhdr}
\fancypagestyle{plain}{%  the preset of fancyhdr 
\fancyhf{} % clear all header and footer fields
\fancyfoot[R]{
\includegraphics[width=2cm]{LogoFCFMBUAP (1).png}}
\fancyfoot[L]{{\bfseries{\thedate{}}} a las {\bfseries{\mycurrenttime{}}} horas{} (GMT-6, H. Puebla de Zaragoza, Pue)}
\fancyhead[L]{Electrodinámica}
\fancyhead[R]{\theauthor}
}
\makeatletter
\def\@maketitle{%
\newpage
\null
\vskip 1em%
\begin{center}%
\let \footnote \thanks
{\LARGE \@title \par}%
\vskip 1em%
%{\large \@date}%
\end{center}%
\par
\vskip 1em}
\makeatother

\usepackage{lipsum}  
\usepackage{cmbright}


%%%%%%%%%%%%%%%%%%%%%%%%%%%%%%%%%%%%%%%%%%%%%%%%%%%%%%%%%%%%%%%%%%%%%%%%%%%%%%%%%%%%%%%%%%%%%%%%%%%%%%%%%%%%%%%%%%%%%%%%%%%%%%%%%%%%%%%%%%%%%%%%%%%%%%%%%%%%%%%%%%%%%%%%%%%%%%%%%%%%%%%%%%%%%%%%%%%%%%

\begin{document}



\maketitle


%%%%%%%%%%%%%%%%%%% Datos del profesor, campus y aula

\begin{mybox}
\noindent\begin{tabular}{@{}ll}
	{\bfseries{Profesor}} & Dr. Iván Fuentecilla Cárcamo\\
	\textcolor{prussianblue}{\bfseries{Campus}} / \textcolor{trueblue}{\bfseries{Facultad}}  & \textcolor{prussianblue}{\bfseries{C.U. BUAP}} /  \textcolor{trueblue}{\bfseries{FCFM}} \\
	\textcolor{prussianblue}{\bfseries{Edificio}} /  \textcolor{trueblue}{\bfseries{Salón}}    & 	\textcolor{prussianblue}{\bfseries{EMA4}} /  \textcolor{trueblue}{\bfseries{409}}   
\end{tabular} 
\end{mybox} \vspace{1cm}




	





\begin{MyColorPar}{orange(colorwheel)}

\end{MyColorPar}


\tableofcontents \newpage

\begin{center}
\begin{tikzpicture}
	% Dibuja círculos rellenos en diferentes colores
	\fill[bittersweet] (0,0) circle (0.17cm);
	\fill[Apple Green] (0.6,0) circle (0.17cm);
	\fill[trueblue] (1.2,0) circle (0.17cm);
	\fill[Aguamarina] (1.8,0) circle (0.17cm);
	\fill[Sun] (2.4,0) circle (0.17cm);
	\fill[orange(colorwheel)] (3.0,0) circle (0.17cm);
	\fill[Pesto] (3.6,0) circle (0.17cm);
	\fill[blue-violet] (4.2,0) circle (0.17cm);
	\fill[ticklemepink] (4.8,0) circle (0.17cm);
	\fill[Gallery] (5.4,0) circle (0.17cm);
	\fill[Mantle] (6.0,0) circle (0.17cm);
	\fill[Nevada] (6.6,0) circle (0.17cm);
	\fill[onyx] (7.2,0) circle (0.17cm);
	% Agrega más círculos o modifica los colores según sea necesario
\end{tikzpicture}
\end{center}  \vspace{2cm}




\section{Unit 2: Electric fields in dielectric media}

\begin{figure}[H]
    \centering
	    \includegraphics[width=0.7\textwidth]{1.png} % ajusta el ancho según convenga
    \caption{Descripción clara y concisa de la figura.}
    \label{fig:mi_figura2}
\end{figure}


Potencial:

\[U = - \vec{P} \cdot \vec{E} \]

Potencial asociado a una polarización. 

\begin{figure}[H]
    \centering
	    \includegraphics[width=0.7\textwidth]{2.png} % ajusta el ancho según convenga
    \caption{Descripción clara y concisa de la figura.}
    \label{fig:mi_figura2}
\end{figure}

Caso extremo cuando tenemos un campo eléctrico bajo, temperatura alta, hace que haya una alineación muy random.

\textbf{Polarización aosciada a un elemento de volumen}

\[ \vec{P} = \frac{1}{\Delta{v}} \sum\limits_{i=!}^{n}\vec{P}_{m}  \]

Como están dispuestos aleatoriamente lo anterior obtiene un valor de $0$. 

\textbf{Polarización en esos casos} es cercano a cero o es cero.

Cuando hablamos de algo que es aleatoria tenemos distribuciones de probabilidad, el problema es tremendamente grande.

\[ \psi_{i} \quad E_{i} \quad \left. P(E_{i}) \right|_{T}  \]

(Mecánica estadística)

\textbf{Boltzmann distribution}

Pasamos ahora a la parte forma, a la parte de cálculo:

El objetivo final de esto anterior es:

\textbf{Plarizabilidad}

\[ \vec{P}_{m} = \alpha \vec{E}_{m} \]

\[ \alpha(R_{0}) \]


La polarizabilidad es una cantidad molecular.

El objetivo de ahora es observar la dependencia de la \textbf{polarizabilidad}, la temperatura debe de tener un papel muy importante.


Podemos encontrar una relación entre la \textbf{constante dieléctrica}, la \textbf{ polarizabilidad} y la \textbf{temperatura}. 

\section{Polar molecules in electric field}

Ahora las moléculas tienen un momento dipolar permanente. 

Tenemos la polarización $P$ dada como:


\[ \vec{P} = \frac{1}{\Delta{V}} \sum\limits_{m} \vec{P}_{m}  \]

Esto anterior puede tener un valor límite que se le denomina de saturación, el cual recibe el siguiente tratamiento o resultado:


Llegar a ese valor es ideal.

\[ \vec{P}_{s} \quad \to \quad \text{Saturation polarization}  \]


Vamos a resolver la siguiente integral.


Tenemos una energía potencial asociado a ese dipolo inmerso en el campo eléctrico.

Para cada molécula, vamos a tener un momento dipolar. Entonces el campo eléctrico anterior es microscópico.


\begin{figure}[H]
    \centering
	    \includegraphics[width=0.7\textwidth]{3.png} % ajusta el ancho según convenga
    \caption{Descripción clara y concisa de la figura.}
    \label{fig:mi_figura2}
\end{figure}

El ángulo $\theta$ se forma entre el vector de \textbf{momento dipolar} y el \textbf{campo eléctrico}. 


El momento dipolar tendrá las componentes de:

\[ x,y,z \]

Tenemos el elemento de aleatoriedad, estocasticidad, entonces tenemos que considerar que no será exacto, entonces tendremos diferentes ángulos.


\[ \underbrace{\langle P_{0} \cos \theta \rangle}_{\text{valor esperado}} =  \int\limits  \]



\textbf{Boltzmann distribution}

\[ \alpha(T) \quad \Bigg{|} \quad  P(E_{i}) \alpha e^{ -\frac{E}{K_{B}T} }  \quad K_{B}=\text{Boltzmann constant}  \]


Tenemos así que el \textbf{valor esperado} es la siguiente integral considerando el \textbf{ángulo sólido} $\Omega$


\[ \langle P_{0}\cos\theta \rangle = \frac { \int\limits P_{0}\cos\theta e^{ \frac{ P_{0}E_{m}\cos\theta }{K_{B}T}  } d\Omega  }{ \int\limits e^{\frac{P_{0}E_{m}\cos\theta }{K_{B}T}  d\Omega }  } \]

Tenemos el siguiente resultado:


\begin{figure}[H]
    \centering
	    \includegraphics[width=0.7\textwidth]{4.png} % ajusta el ancho según convenga
    \caption{Descripción clara y concisa de la figura.}
    \label{fig:mi_figura2}
\end{figure}

\begin{figure}[H]
    \centering
	    \includegraphics[width=0.7\textwidth]{5.png} % ajusta el ancho según convenga
    \caption{Descripción clara y concisa de la figura.}
    \label{fig:mi_figura2}
\end{figure}

Realizamos la siguiente integración, tenemos:


Tenemos: 

\[ \frac{ 2\pi P_{0}  \int\limits_{0}^{\pi} \cos \theta e^{y\cos \theta}\sen\theta d\theta  }{ 2\pi \int\limits e^{y\cos\theta}\sen\theta d\theta  }  \]


Tenemos:


\[\left[ \int\limits_{0}^{\pi}  e^{y\cos\theta}\sen\theta d\theta = -\frac{1}{y}e^{y\cos\theta} \right]_{0}^{\pi} = -\frac{1}{y} \left[ e^{-y}   -e^{y}\right]  \]

Con los cambios de variable:


\[ u=\cos\theta \quad du=-\sen\theta d\theta  \]

\[ \sen{hx} = \frac{ e^{x}-e^{-x} }{2} \]

\[ \cos{hx} = \frac{e^{x}+e^{-x}}{2} \]

La segunda integral es:


\[ \int\limits_{0}^{\pi}  \cos\theta e^{y\cos\theta}\sen\theta d\theta  = \left. - \cos\theta \frac{e^{y\cos\theta}}{y} + \frac{e^{y\cos\theta}}{y^{2}} \right]^{\pi}_{0}  \]

Con los cambios de variable:


\[ u=\cos\theta, du = -\sen\theta d\theta  \]

\[\int\limits \cos\theta e^{y\cos\theta}\sen\theta d\theta = - \int\limits u e^{yu}du \]


\[ \frac{1}{y} e^{yu}du = v \quad (v^{\prime}=e^{yu})= - \left[ \frac{ue^{yu}}{y} - \int\limits \frac{e^{yu}}{y}du  \right]   \]


Tenemos:


\[ = - \frac{e^{yu}}{y} + \frac{1}{y^{2}} e^{yu} = - \frac{ \cos\theta e^{y\cos\theta} }{y} + \frac{e^{y\cos\theta}}{y^{2}}  \]




\[ \int\limits_{0}^{\pi}  \cos\theta e^{y\cos\theta}\sen\theta d\theta  = \left. - \cos\theta \frac{e^{y\cos\theta}}{y} + \frac{e^{y\cos\theta}}{y^{2}} \right]^{\pi}_{0} = \frac{e^{-y}}{y} + \frac{e^{-y}}{y^{2}} + \frac{e^{y}}{y}  - \frac{e^{y}}{y^{2}}\]

\[ \frac{2}{y}\cos{hy} - \frac{2}{y^{2}}\sen{hy}  \]

\[ \langle P_{0}\cos\theta \rangle = P_{0} \left[ \frac{ \cos{hy} }{\sen{hy}} - \frac{1}{y} \right] \]

\[ \boxed{\langle P_{0}\cos\theta  \rangle}  = P_{0} \left[ cot{hy} - \frac{1}{y}  \right] \]

Ecuación 5.21


Tarea problema 5.4.

\textbf{Activity}: Using the expansion:

\[\cot{hx} = \frac{1}{x} + \frac{x}{3} + \ldots +  \]

At the smallets order, obtain from \textbf{Langevin equation}

\[  \langle P_{0} \cos\theta  \rangle = \frac{P_{0}^{2}E_{m} }{3K_{B}T} \]


\[ \boxed{ \langle P_{0}\cos\theta \rangle  } = \frac{ P_{0}^{2} }{3} \frac{E_{m}}{K_{B}T} \]

\begin{figure}[H]
    \centering
	    \includegraphics[width=0.7\textwidth]{6.png} % ajusta el ancho según convenga
    \caption{Descripción clara y concisa de la figura.}
    \label{fig:mi_figura2}
\end{figure}


\[
\textbf{Tenemos:} \qquad 
\langle P_{0}\cos\theta\rangle = \frac{P_{0}^{2}}{3}\frac{E_{m}}{k_{B}T}
\]

Sabemos además que la polarización inducida (por partícula) puede escribirse como:

\[
P_{m} = \alpha E_{m}
\]

Comparando ambas expresiones, observamos que el término que acompaña a \(E_{m}\) 
en la primera ecuación corresponde directamente a la \textbf{polarizabilidad}:

\[
\boxed{ \alpha(T) = \frac{P_{0}^{2}}{3k_{B}T} }
\]

\vspace{0.5cm}

Si existen \(N\) dipolos por unidad de volumen, la \textbf{polarización macroscópica} es:

\[
\vec{P} = N \langle P_{0}\cos\theta\rangle = N \alpha(T) \vec{E}_{m}
\]

Finalmente, usando la relación 
\(\vec{P} = \varepsilon_{0}\chi_{e}\vec{E}_{m}\),
se obtiene la \textbf{susceptibilidad eléctrica}:

\[
\boxed{ \chi_{e}(T) = \frac{N P_{0}^{2}}{3\varepsilon_{0}k_{B}T} }
\]



 
\end{document}